%% Question bank chapter file - generated by bank_tools.py build (do not edit while the bank is still being built:
%% it is regenerated from the per-test data in Question Bank\data). Every question is one bankq block:
%%   \begin{bankq}[id=..., source={...}, paper=.., qno=.., type=..., subtopic={...}, answer={...}, answer_src=key|solved,
%%                 working={...}, difficulty=1-10, diff_parts={K S I C F}, diff_note={...},
%%                 group=<passage id or empty>, also_in={duplicate copies in other tests}]
%% difficulty = min(10, K+S+I+C+F-1): K concepts 1-3, S steps 1-4, I insight 0-2, C computation 0-2, F format 0-1
%%   ...question LaTeX, options, figures (\bfig / \bfigi -> img/)...
%%   \end{bankq}
%% Passages shared by several questions are bankpassage blocks; their questions carry group=<passage id>.
\documentclass[10pt]{article}
\usepackage[a4paper,margin=1.8cm]{geometry}
\usepackage[T1]{fontenc}
\usepackage{amsmath,amssymb,mathtools}
\usepackage{graphicx}
\usepackage[export]{adjustbox}
\usepackage{xcolor,needspace,keyval}
\usepackage{enumitem}
\usepackage{booktabs,array,tabularx}
\usepackage[hidelinks]{hyperref}
\setlength{\parindent}{0pt}
\setlength{\parskip}{3pt}
\makeatletter
\define@key{bq}{id}{\def\bq@id{#1}}\define@key{bq}{source}{\def\bq@src{#1}}\define@key{bq}{paper}{\def\bq@paper{#1}}
\define@key{bq}{qno}{\def\bq@qno{#1}}\define@key{bq}{type}{\def\bq@type{#1}}\define@key{bq}{subtopic}{\def\bq@sub{#1}}
\define@key{bq}{answer}{\def\bq@ans{#1}}\define@key{bq}{answer_src}{\def\bq@asrc{#1}}\define@key{bq}{working}{\def\bq@work{#1}}
\define@key{bq}{group}{\def\bq@group{#1}}\define@key{bq}{also_in}{\def\bq@also{#1}}
\define@key{bq}{difficulty}{\def\bq@diff{#1}}\define@key{bq}{diff_parts}{\def\bq@dparts{#1}}\define@key{bq}{diff_note}{\def\bq@dnote{#1}}
\newenvironment{bankq}[1][]{\def\bq@work{}\def\bq@also{}\def\bq@group{}\def\bq@diff{?}\def\bq@dparts{}\def\bq@dnote{}\setkeys{bq}{#1}%
  \par\needspace{4\baselineskip}\medskip\noindent{\small\textbf{\bq@id}\hfill\color{gray}[\bq@src, Paper \bq@paper, Q\bq@qno]
  \bq@type\ $\cdot$ \bq@sub\ $\cdot$ \textbf{D\bq@diff}}\par\nopagebreak\relax}
  {\par{\small\color{teal}Answer: \textbf{\bq@ans}\ifx\bq@work\@empty\else\ (solved: \bq@work)\fi
  \ifx\bq@also\@empty\else\par\color{gray}Also in: \bq@also\fi}\par}
\define@key{bp}{id}{\def\bp@id{#1}}\define@key{bp}{source}{\def\bp@src{#1}}\define@key{bp}{questions}{\def\bp@qs{#1}}
\newenvironment{bankpassage}[1][]{\setkeys{bp}{#1}\par\needspace{4\baselineskip}\bigskip\noindent
  {\small\textbf{Passage \bp@id}\hfill\color{gray}[\bp@src; questions \bp@qs]}\par\nopagebreak\relax\itshape}{\par}
\makeatother
\newcommand{\bfig}[2]{\par\smallskip{\centering\includegraphics[width=#1pt,max width=0.62\linewidth,max height=0.32\textheight]{img/#2}\par}\smallskip}
\newcommand{\bfigi}[2]{\includegraphics[width=#1pt,max width=\linewidth,max height=0.32\textheight,valign=c]{img/#2}}
\newcommand{\bmissing}[1]{\textit{[figure missing in the source conversion: \detokenize{#1}]}}
\newenvironment{qoptions}{\begin{itemize}[leftmargin=1.8em,itemsep=1pt,topsep=2pt,label={}]}{\end{itemize}}
\newcommand{\opt}[2]{\item[(#1)] #2}
\providecommand{\fwstart}{}\providecommand{\fwend}{}\providecommand{\qfullsplit}{}
\newcolumntype{L}[1]{>{\raggedright\arraybackslash}p{#1}}
\newcolumntype{C}[1]{>{\centering\arraybackslash}p{#1}}
\newcolumntype{R}[1]{>{\raggedleft\arraybackslash}p{#1}}
\providecommand{\sectionhdr}[1]{}\providecommand{\subsectionhdr}[1]{}


\begin{document}
{\Large\bfseries POC --- Question Bank}\par\smallskip
{\small Chemistry chapter C24. Questions from all converted test papers, grouped by source series. Answers: printed key where it matches the paper, otherwise solved (answer\_src=solved, with working). Exact repeats across papers appear once, with also\_in listing the other copies.}\par\medskip\hrule

\section*{TOAS papers}
\begin{bankq}[id=TOAS08-P2-CHEM-S2-Q10, source={TOAS - Test 08 (15 May 2022)}, paper=2, qno=10, type=single, subtopic={Qualitative tests}, answer={A}, answer_src=key, difficulty=1, diff_parts={K1 S1 I0 C0 F0}, diff_note={concepts: qualitative analysis lists | steps: reactions of each substance; match}]
Match the list-
\begin{center}\small\renewcommand{\arraystretch}{1.1}\setlength{\extrarowheight}{3.5pt}\setlength{\tabcolsep}{3pt}
\begin{tabular}{|c|c|c|l|}
\hline
\multicolumn{2}{|c|}{\parbox[t]{3.6cm}{\centering List-I\\ (Pair of compound)}} & \multicolumn{2}{c|}{\parbox[t]{2.4cm}{\centering List-II\\ (Differentiated by)}}\\[3pt]
\hline
(P) & \bfigi{132.0}{TOAS08_p2_chem_s2_q10_1.png} & (1) & Fehling test\\[3pt]
(Q) & \bfigi{61.8}{TOAS08_p2_chem_s2_q10_2.png} & (2) & Tollen's test\\[3pt]
(R) & \bfigi{76.9}{TOAS08_p2_chem_s2_q10_3.png} & (3) & Na-metal\\[3pt]
(S) & \bfigi{104.3}{TOAS08_p2_chem_s2_q10_4.png} & (4) & $\mathrm{NaHCO_3}$\\[3pt]
 & & (5) & Cold $\mathrm{KMnO_4}$\\[3pt]
\hline
\end{tabular}
\end{center}
\begin{qoptions}
\opt{A}{$P\to1,2,;Q\to2;R\to1,2,4;S\to1,2,5$}
\opt{B}{$P\to1,2;Q\to1,2;R\to1,2;S\to1,2$}
\opt{C}{$P\to1;Q\to1,2;R\to1,2,4;S\to1,2,5$}
\opt{D}{$P\to1,2;Q\to2;R\to1,4;S\to5$}
\end{qoptions}
\end{bankq}
\begin{bankq}[id=TOAS08-P2-CHEM-S3-Q06, source={TOAS - Test 08 (15 May 2022)}, paper=2, qno=6, type=numerical, subtopic={Qualitative tests}, answer={4.00}, answer_src=key, difficulty=2, diff_parts={K1 S1 I0 C0 F1}, diff_note={concepts: Tollens' test | steps: identify aldehydes and alpha-hydroxy ketones; count positive tests}]
Which of the following can give silver mirror with Tollen's reagent :
\bfig{245.2}{TOAS08_p2_chem_s3_q06_1.png}
\end{bankq}
\begin{bankq}[id=TOAS09-P1-CHEM-S4-Q09, source={TOAS - Test 09 (30 May 2022)}, paper=1, qno=9, type=numerical, subtopic={Qualitative tests}, answer={6}, answer_src=key, difficulty=3, diff_parts={K1 S2 I0 C0 F1}, diff_note={concepts: Tollens' test | steps: identify aldehydes, formic acid and terminal alkynes; positive tests; count}]
How many of the following will give silver mirror with Tollen's reagent ?\par
(i) Ph--CH=O, (ii) HCHO, (v) HCOOH, (vi) MeCOOH, (viii) HC$\equiv$CH, and the structures below, in order (iii), (iv), (vii), (ix), (x):
\bfig{166.8}{TOAS09_p1_chem_s4_q09_1.png}
\end{bankq}
\begin{bankq}[id=TOAS09-P2-CHEM-S2-Q06, source={TOAS - Test 09 (30 May 2022)}, paper=2, qno=6, type=numerical, subtopic={Qualitative tests}, answer={3}, answer_src=solved, working={Positive Fehling: HCHO, HCOOH and cyclohexanecarbaldehyde (aromatic aldehyde, ketone, acetal, hemiketal, diol do not)}, difficulty=3, diff_parts={K1 S2 I0 C0 F1}, diff_note={concepts: Fehling test | steps: aliphatic aldehydes and formic acid give a positive test; aromatic aldehydes do not; count}]
How many of the following will give positive Fehling test :
\bfig{134.0}{TOAS09_p2_chem_s2_q06_1.png}
\end{bankq}
\begin{bankq}[id=TOAS09-P2-CHEM-S2-Q07, source={TOAS - Test 09 (30 May 2022)}, paper=2, qno=7, type=numerical, subtopic={Qualitative tests}, answer={6}, answer_src=solved, working={Not iodoform: iPrOTs, tert-amyl alcohol, pentan-3-one, acetyl chloride, acetic anhydride, propanal (EtOH and pyruvic acid are positive)}, difficulty=3, diff_parts={K1 S2 I0 C0 F1}, diff_note={concepts: iodoform test | steps: methyl ketones and CH3CH(OH)R give the test; examine each compound; count the negatives}]
How many of the following will not give iodoform test :
\bfig{149.4}{TOAS09_p2_chem_s2_q07_1.png}
\end{bankq}
\begin{bankq}[id=TOAS10-P1-CHEM-S3-Q07, source={TOAS - Test 10 (12 Jun 2022)}, paper=1, qno=7, type=numerical, subtopic={Qualitative tests}, answer={3.00}, answer_src=key, difficulty=2, diff_parts={K1 S1 I0 C0 F1}, diff_note={concepts: reactions and tests | steps: check each of the four statements or reactions; count correct ones}]
Find number of reaction correctly matched, with their major product (or statement written with them) :
\bfig{163.1}{TOAS10_p1_chem_s3_q07_1.png}
(4) CH$_3$CHO gives positive Tollen's test as well as Lucas test
\bfig{157.3}{TOAS10_p1_chem_s3_q07_2.png}
\bfig{164.1}{TOAS10_p1_chem_s3_q07_3.png}
\bfig{154.2}{TOAS10_p1_chem_s3_q07_4.png}
\end{bankq}
\begin{bankq}[id=TOAS11-P1-CHEM-S2-Q08, source={TOAS - Test 11 (27 Jun 2022)}, paper=1, qno=8, type=multi, subtopic={Qualitative tests}, answer={A,B,D}, answer_src=key, difficulty=2, diff_parts={K1 S1 I0 C0 F1}, diff_note={concepts: Fehling and Tollens tests | steps: identify the functional groups of each pair; test each reagent}]
Among the following pair of compounds, identify the pair(s) which can be differentiated by Fehling as well as Tollen's reagent ?
\begin{qoptions}
\opt{A}{\bfigi{74.9}{TOAS11_p1_chem_s2_q08_A.png}}
\opt{B}{\bfigi{102.3}{TOAS11_p1_chem_s2_q08_B.png}}
\opt{C}{\bfigi{107.0}{TOAS11_p1_chem_s2_q08_C.png}}
\opt{D}{\bfigi{100.9}{TOAS11_p1_chem_s2_q08_D.png}}
\end{qoptions}
\end{bankq}
\begin{bankq}[id=TOAS11-P1-CHEM-S3-Q04, source={TOAS - Test 11 (27 Jun 2022)}, paper=1, qno=4, type=numerical, subtopic={Qualitative tests}, answer={3.00}, answer_src=key, difficulty=3, diff_parts={K1 S2 I0 C0 F1}, diff_note={concepts: tests for a carbonyl compound | steps: identify X; 2,4-DNP; Tollens and Fehling; iodoform and bisulphite tests}]
\bfig{174.1}{TOAS11_p1_chem_s3_q04_1.png}
How many statement is/are correct for [X]\par
(a) It can give orange colour compound with 2,4--D.N.P.\par
(b) It can give positive Silver mirror test.\par
(c) It can be oxidised by Feheling solution\par
(d) It can give yellow ppt. with NaOI,$\Delta$\par
(e) It can be identified by NaHSO$_3$
\end{bankq}
\begin{bankq}[id=TOAS12-P1-CHEM-S1-Q06, source={TOAS - Test 12 (10 Jul 2022)}, paper=1, qno=6, type=multi, subtopic={Qualitative tests}, answer={A,B,C}, answer_src=key, difficulty=2, diff_parts={K1 S1 I0 C0 F1}, diff_note={concepts: distinguishing carbonyl derivatives | steps: identify the products P and Q; test each reagent on both}]
\bfig{133.4}{TOAS12_p1_chem_s1_q06_1.png}
Products formed by P \& Q can not be differentiated by:
\begin{qoptions}
\opt{A}{2, 4 DNP}
\opt{B}{Lucas reagent (ZnCl$_2$) conc. HCl}
\opt{C}{NaHSO$_3$}
\opt{D}{Fehlings solution}
\end{qoptions}
\end{bankq}
\begin{bankq}[id=TOAS12-P1-CHEM-S2-Q08, source={TOAS - Test 12 (10 Jul 2022)}, paper=1, qno=8, type=numerical, subtopic={Qualitative tests}, answer={4.00}, answer_src=key, difficulty=3, diff_parts={K1 S2 I0 C0 F1}, diff_note={concepts: Fehling test | steps: identify aldehydes; exclude ketones, aromatic aldehydes and hemiacetals; count}]
How many compounds will give positive Fehling test :\par
(i) CH$_3$CHO\quad (ii) PhCHO\quad (iii) CH$_3-$CO$-$CH$_3$\quad (iv) CH$_3-$CH(OH)(OCH$_3$)
\bfig{223.9}{TOAS12_p1_chem_s2_q08_2.png}
\end{bankq}
\begin{bankq}[id=TOAS12-P2-CHEM-S3-Q04, source={TOAS - Test 12 (10 Jul 2022)}, paper=2, qno=4, type=numerical, subtopic={Qualitative tests}, answer={64.00}, answer_src=key, difficulty=3, diff_parts={K1 S2 I0 C0 F1}, diff_note={concepts: tests for compounds | steps: Tollens test positives; iodoform positives; two-digit answer}]
\bfig{227.5}{TOAS12_p2_chem_s3_q03_4.png}
(j) CH$_3$--CH(OH)--Ph\par
Answer in two digit form XY\par
X = Number of compounds giving (+) test with AgNO$_3$ / NH$_4$OH\par
Y = Number of compounds giving yellow colour with I$_2$+ NaOH
\end{bankq}
\begin{bankq}[id=TOAS13-P1-CHEM-S3-Q05, source={TOAS - Test 13 (24 Jul 2022)}, paper=1, qno=5, type=numerical, subtopic={Qualitative tests}, answer={5.00}, answer_src=key, difficulty=3, diff_parts={K1 S2 I0 C0 F1}, diff_note={concepts: iodoform test | steps: find compounds with CH3CO or CH3CH(OH); check each of the eight; count}]
How many compounds gives iodoform test with NaOH $+$ I$_2$ :
\bfig{258.7}{TOAS13_p1_chem_s3_q05_1.png}
\end{bankq}
\begin{bankq}[id=TOAS15-P2-CHEM-S2-Q05, source={TOAS - Test 15 (22 Aug 2022)}, paper=2, qno=5, type=numerical, subtopic={Qualitative tests}, answer={3.00}, answer_src=key, difficulty=3, diff_parts={K1 S2 I0 C0 F1}, diff_note={concepts: aldehyde tests with Fehling and DNP reagents | steps: aldehydes respond to both; ketones only to DNP; count the aldehydes among the structures}]
How many compound show Fehling and 2, 4 DNP test.\par
\noindent\bfigi{56.3}{TOAS15_p2_chem_s2_q05_2.png}\hfill\bfigi{27.9}{TOAS15_p2_chem_s2_q05_1.png}\hfill\bfigi{43.2}{TOAS15_p2_chem_s2_q05_3.png}\par
\noindent\bfigi{84.5}{TOAS15_p2_chem_s2_q05_5.png}\hfill\bfigi{48.7}{TOAS15_p2_chem_s2_q05_4.png}\par
\noindent\bfigi{58.8}{TOAS15_p2_chem_s2_q05_7.png}\hfill\bfigi{51.9}{TOAS15_p2_chem_s2_q05_6.png}
\end{bankq}
\begin{bankq}[id=TOAS17-P2-CHEM-S2-Q06, source={TOAS - Test 17 (18 Sep 2022)}, paper=2, qno=6, type=single, subtopic={Qualitative tests}, answer={6}, answer_src=key, difficulty=2, diff_parts={K1 S1 I1 C0 F0}, diff_note={concepts: Tollens and Fehling tests | steps: aliphatic aldehydes and formic acid respond to both; aromatic aldehydes and ketones do not}]
Identify total number of compounds which give both tollen's and fehling test ?
\par\noindent(i) HCHO\quad(ii) CH$_3$--CHO\quad(iii) Ph--CHO\quad(iv) CH$_3$--C$\equiv$C--CH$_3$\par
\noindent(v) \bfigi{43.6}{TOAS17_p2_chem_s2_q06_2.png}\quad(vi) Ph--CH$_2$--CHO\par
\noindent(vii) HCOOH\quad(viii) CH$_3$--CO--CH$_2$--OH\par
\noindent(ix) CH$_3$--COOH\quad(x) Ph--CO--Ph
\end{bankq}
\begin{bankq}[id=TOAS18-P2-CHEM-S2-Q10, source={TOAS - Test 18 (30 Oct 2022)}, paper=2, qno=10, type=multi, subtopic={Qualitative tests}, answer={A,B,C,D}, answer_src=key, difficulty=2, diff_parts={K1 S1 I0 C0 F1}, diff_note={concepts: haloform test products | steps: identify methyl ketone or methyl carbinol in each product}]
Reaction in which product can show positive Haloform test :
\begin{qoptions}
\opt{A}{\bfigi{79.3}{TOAS18_p2_chem_s2_q10_A.png}}
\opt{B}{\bfigi{138.9}{TOAS18_p2_chem_s2_q10_B.png}}
\opt{C}{\bfigi{123.9}{TOAS18_p2_chem_s2_q10_C.png}}
\opt{D}{\bfigi{92.5}{TOAS18_p2_chem_s2_q10_D.png}}
\end{qoptions}
\end{bankq}
\begin{bankq}[id=TOAS19-P1-CHEM-S3-Q05, source={TOAS - Test 19 (05 Feb 2023)}, paper=1, qno=5, type=numerical, subtopic={Practical organic chemistry}, answer={34.04}, answer_src=key, difficulty=3, diff_parts={K1 S1 I0 C1 F1}, diff_note={concepts: Carius method | steps: moles of bromine in the silver bromide; percentage}]
In carius method of estimation of halogen 0.15 g of an organic compound gave 0.12 g of AgBr. Find out the percentage of bromine in the compound?
\end{bankq}
\begin{bankq}[id=TOAS19-P1-CHEM-S3-Q06, source={TOAS - Test 19 (05 Feb 2023)}, paper=1, qno=6, type=numerical, subtopic={Practical organic chemistry}, answer={21.95}, answer_src=key, difficulty=3, diff_parts={K1 S1 I0 C1 F1}, diff_note={concepts: combustion analysis | steps: mass of carbon from carbon dioxide; percentage}]
On complete combustion, 0.246 g of an organic compound gave 0.198 g of carbon dioxide and 0.1014g of water. Determine the percentage composition of carbon in the compound
\end{bankq}
\end{document}
